\section{Jacobi scattering and theta functions}
\label{sec:jacobi-theta}

\subsection{The twisted volume line}

Fix a root chamber and an integral basis $e_1,e_2$ of $\widetilde M$.  Put
\begin{equation}\label{eq:primitive-form}
 \kappa=e_1+e_2,
 \qquad \Omega_0=d\log z_1\wedge d\log z_2,
 \qquad \Omega_\kappa=z^\kappa\Omega_0.
\end{equation}
Let $J(q_1,q_2)=(q_2,-q_1)$ and denote the invariant logarithmic derivation
in direction $v$ by $\partial_v$.

\begin{proposition}[Shifted Hamiltonian bracket]
\label{prop:shifted-bracket}
The inverse bivector $\pi_\kappa=z^{-\kappa}\pi_0$ of $\Omega_\kappa$ has
Hamiltonian vector fields
\begin{equation}\label{eq:shifted-Hamiltonian}
 X_q^\kappa=z^{q-\kappa}\partial_{Jq},
\end{equation}
and
\begin{equation}\label{eq:shifted-bracket}
 [X_q^\kappa,X_r^\kappa]
 =-\det(q,r)X_{q+r-\kappa}^\kappa.
\end{equation}
Thus, after writing $q=k+\kappa$ and $r=\ell+\kappa$, the output is indexed
by $k+\ell$.
\end{proposition}

\begin{proof}
The Hamiltonian field of $z^q$ for $z^{-\kappa}\pi_0$ is
\eqref{eq:shifted-Hamiltonian}.  The bracket of two monomial derivations has
direction
\[
 \langle Jq,r-\kappa\rangle Jr
 -\langle Jr,q-\kappa\rangle Jq
 =-\det(q,r)J(q+r-\kappa),
\]
which gives \eqref{eq:shifted-bracket}.
\end{proof}

The lifted slab map at $V$ is symplectic for $\Omega_0$, but
\begin{equation}\label{eq:slab-conformal}
 \widetilde\theta_V^*\Omega_\kappa
 =g_V^{\langle N_V,\kappa\rangle}\Omega_\kappa,
 \qquad
 \bigl(\langle N_V,\kappa\rangle\bigr)_V=(-2,-2,2,-2).
\end{equation}
Choose a local line frame $e$ and identify frames across the slab by
\begin{equation}\label{eq:line-transition}
 \widetilde\theta_V^*e'=g_V^{-\langle N_V,\kappa\rangle}e.
\end{equation}
The products $\Omega_\kappa\otimes e$ then glue.  We denote the resulting
line by $\Lcal_\kappa$.

The transition law therefore requires a line bundle.

\begin{proposition}[No scalar frame]\label{prop:no-scalar-frame}
The scalar $\kappa$-twisted Hamiltonian algebra containing the two singleton
fields is not preserved under conjugation by any lifted slab transition.
Consequently its local brackets do not glue to a bracket on the structure
sheaf.
\end{proposition}

\begin{proof}
If $p_V$ is the slab exponent and
\[
 S_V=\frac1{d_V}\log(1+z^{p_V})\partial_{N_V},
\]
then for $k=q-\kappa$ one computes
\begin{equation}\label{eq:mixed-slab-bracket}
 [S_V,X_q^\kappa]
 =\sum_{r\geq1}\frac{(-1)^{r+1}}{d_V}z^{k+rp_V}
 \left(
  \frac{\langle N_V,k\rangle}{r}\partial_{Jq}
  -\langle Jq,p_V\rangle\partial_{N_V}
 \right).
\end{equation}
With respect to $\Omega_\kappa$, the divergence of the $r$th summand has
coefficient
\[
 \frac{(-1)^r}{d_V}
 \langle Jq,p_V\rangle\langle N_V,\kappa\rangle.
\]
For $r=1$ and the singleton exponents $q_1=(2,1)$ and $q_2=(1,2)$, the
eight coefficients are
\[
\begin{array}{c|rrrr}
 &A&B_0&B_1&C\\ \hline
 q_1&2&3&1&-10\\
 q_2&10&3&-7&-2.
\end{array}
\]
Thus every singleton--slab interaction has nonzero divergence.
\end{proof}

\begin{proposition}[Jacobi transport]\label{prop:Jacobi-transport}
The line $\Lcal_\kappa$ on the pathwise pro-root object carries a canonical
Jacobi bracket whose expression in the distinguished chamber is the
$\kappa$-twisted Poisson bracket.  Every lifted slab transition is a Jacobi
isomorphism.  The construction is compatible with finite root restriction
and descends with the pro-Kummer coefficient object.
\end{proposition}

\begin{proof}
In a frame $e$, write a Jacobi bracket as
\[
 \{fe,he\}=
 \bigl(\Lambda(df,dh)+fE(h)-hE(f)\bigr)e.
\]
Under $e'=ae$ the local pair transforms by
\begin{equation}\label{eq:Jacobi-frame-change}
 \Lambda'=a\Lambda,
 \qquad E'=aE+\Lambda^\sharp(da).
\end{equation}
Start at the root chamber with $(\Lambda,E)=(\pi_\kappa,0)$ and transport
the pair across each edge using \eqref{eq:line-transition} and
\eqref{eq:Jacobi-frame-change}.  Every vertex of the path tree has a unique
route from the root, so the construction is unambiguous and compatible with
restriction to finite subtrees.  Change of frame and pullback preserve the
Jacobi identity.  The deck action transports the same data, giving descent.
\end{proof}

Equivalently, $\Lcal_\kappa^\vee\setminus0$ carries the Poissonization of
this Jacobi structure.  The line-bundle formulation displays the slab
multiplier in \eqref{eq:line-transition} explicitly.

\begin{definition}[Jacobi walls and consistency]\label{def:Jacobi-wall}
Fix a finite root-stack presentation and reduce the coefficient ring modulo $I^N$.  A
\emph{Jacobi wall} is a pair $(\mathfrak d,\Theta_{\mathfrak d})$, where
$\mathfrak d$ is an oriented rational affine segment, ray, or full chord in
the smooth locus and $\Theta_{\mathfrak d}$ is a locally constant
unipotent Jacobi automorphism along $\mathfrak d$, congruent to the identity
modulo $I$.  The automorphism is read in the local frame of
$\Lcal_\kappa$ and is transported between wall charts by
Proposition~\ref{prop:Jacobi-transport}.

For a transverse path $\gamma$, its \emph{path-ordered product}
$\Theta_\gamma$ is the product of the crossed wall automorphisms in crossing
order, using the inverse when the coorientation is reversed.  A locally
finite collection of Jacobi walls is \emph{consistent at a joint} if the
path-ordered product around a sufficiently small positively oriented loop
about that joint is the identity.  It is \emph{contractibly consistent} if
$\Theta_\gamma$ depends only on the homotopy class of $\gamma$ relative to
its endpoints for homotopies through compact transverse paths.
\end{definition}

All these notions are finite-stage notions.  A wall structure on the
pro-Kummer object means a compatible system of such structures as the root
stage and $N$ vary.  This convention prevents inverse-limit notation from
hiding an infinite path-ordered product.

\subsection{The singleton joint}

Let $t_1,t_2$ be the two singleton parameters, $I=(t_1,t_2)$, and put
\begin{equation}\label{eq:singleton-generators}
 E_{a,b}=t_1^at_2^bX_{\kappa+(a,b)}^\kappa,
 \qquad (a,b)\in\NN^2\setminus\{(0,0)\}.
\end{equation}
Then
\begin{equation}\label{eq:singleton-bracket}
 [E_{a,b},E_{c,d}]
 =-\det\bigl(\kappa+(a,b),\kappa+(c,d)\bigr)E_{a+c,b+d}.
\end{equation}
The two incoming GKT normalizations are
\begin{equation}\label{eq:incoming-singletons}
 \mathsf A=-\frac12E_{1,0},
 \qquad \mathsf B=-\frac12E_{0,1}.
\end{equation}
Since the shifted exponents are $(2,1)$ and $(1,2)$,
\begin{equation}\label{eq:first-bracket}
 [\mathsf A,\mathsf B]= -\frac34E_{1,1}.
\end{equation}

\begin{proposition}[Ordered singleton completion]
\label{prop:singleton-completion}
Let $\widehat{\mathfrak g}_{\mathrm{sing}}$ be the product of the lines
spanned by the $E_{a,b}$, completed in total deformation degree.  The
inverse commutator
\begin{equation}\label{eq:singleton-commutator}
 \bigl(e^{\mathsf A}e^{\mathsf B}e^{-\mathsf A}e^{-\mathsf B}\bigr)^{-1}
\end{equation}
has a unique factorization into ray groups in increasing shifted slope.  The
factorization is finite modulo $I^N$, is compatible under reduction of $N$,
and begins with $\exp(\frac34E_{1,1})$ in the chosen outgoing orientation.
\end{proposition}

\begin{proof}
Formula~\eqref{eq:singleton-bracket} gives an $\NN^2$-grading with finitely
many components modulo $I^N$.  Generators whose shifted exponents lie on one
ray commute.  At each total degree, take the logarithm of the residual word,
decompose its homogeneous part into ray summands, and insert the summands in
slope order.  This is the standard filtered factorization for
two-dimensional scattering diagrams \cite{grosspandharipandesiebert2010}.
Directness of the ray decomposition gives uniqueness.  Equation
\eqref{eq:first-bracket} gives the initial coefficient and the inverse-loop
sign.
\end{proof}

\subsection{Cofactor supports and consistency}

In the plus cofactor chart, the incoming shifted exponents have support
directions
\begin{equation}\label{eq:cofactor-singletons}
 L(2,1)=(-4,3),
 \qquad L(1,2)=(-2,3).
\end{equation}
Each is placed on the full chord through the central point $O$.  Full chords
are forced: a small loop around the endpoint of a truncated degree-one wall
would cross that wall once, while every correction begins in degree two.

If $q=\lambda(2,1)+\mu(1,2)$ with $\lambda,\mu\geq0$, then
\begin{equation}\label{eq:outgoing-cofactor}
 Lq=(-4\lambda-2\mu,3\lambda+3\mu).
\end{equation}
For $\lambda>\mu$ the ray exits one upper facet, for $\mu>\lambda$ it exits
the other, and for $\lambda=\mu$ it runs to the $B$ singularity.  A
correction parallel to an incoming ray is multiplied into the positive half
of the same chord.  Every other correction is placed on its own outgoing
radial support.

\begin{figure}[t]
\centering
\begin{adjustbox}{max width=\textwidth,center}
\begin{tikzpicture}[
  line cap=round,
  line join=round,
  >=Latex,
  every node/.style={font=\scriptsize},
  singletonone/.style={blue!75!black,very thick},
  singletontwo/.style={red!70!black,very thick,densely dashed},
  completion/.style={black!70,thick,-{Latex[length=2mm]}},
  branch/.style={black!45,densely dotted},
  focus/.style={circle,draw=black,fill=white,inner sep=1.8pt,line width=.6pt}
]
% Bounding box.
\draw[white] (-1.35,-1.45) rectangle (8.95,1.75);
% (a) The chart of Delta.
\begin{scope}[x=1.13cm,y=1.13cm]
  \coordinate (pA) at (-1,-1);
  \coordinate (pB) at (-1,1);
  \coordinate (pC) at (3,-1);
  \coordinate (O) at (0,0);
  \fill[black!3] (pA)--(pB)--(pC)--cycle;
  \draw[black!65,thick] (pA)--(pB)--(pC)--cycle;

  \draw[branch] (-.5,-.5)--(1.5,-.5);
  \node[focus,label={[black!70]below left:$A$}] at (-.5,-.5) {};
  \node[focus,label={[black!70]below right:$C$}] at (1.5,-.5) {};

  \draw[singletonone] (-1,.75)--(O)--(1.333,-1);
  \draw[singletontwo] (-.5,.75)--(O)--(.667,-1);
  \fill[blue!75!black] (.667,-.5) circle (1.25pt);
  \fill[red!70!black] (.333,-.5) circle (1.25pt);

  \draw[completion] (O)--(-1,.90);
  \draw[completion] (O)--(-.80,.90);
  \draw[completion] (O)--(-.50,.50);
  \draw[black!35,densely dotted] (-.50,.50)--(pB);
  \node[focus,label={[black!70]right:$B$}] at (-.50,.50) {};

  \fill[black] (O) circle (1.45pt);
  \node[below right=1pt] at (O) {$O$};
  \node[left=1pt] at (-1,.75) {$(-1,\frac34)$};
  \node[above right=0pt] at (-.5,.75) {$(-\frac12,\frac34)$};
  \node[above right=1pt] at (1.333,-1) {$(\frac43,-1)$};
  \node[above left=1pt] at (.667,-1) {$(\frac23,-1)$};
  \node[above left=0pt] at (pB) {$p_B$};
  \node[anchor=south] at (1,1.22)
    {\textbf{(a)} the chart $(s_1,s_2)$ of $\Delta$};
\end{scope}

% Topological annular lift.  Horizontal edges are deleted boundary circles;
% vertical edges are identified.
\begin{scope}[x=1.05cm,y=1.22cm]
  \draw[black!50,densely dashed] (4.35,-1)--(8.35,-1);
  \draw[black!50,densely dashed] (4.35,1)--(8.35,1);
  \draw[black!65] (4.35,-1)--(4.35,1);
  \draw[black!65] (8.35,-1)--(8.35,1);
  \node[left=2pt] at (4.35,0) {$\sim$};
  \node[right=2pt] at (8.35,0) {$\sim$};

  \coordinate (Ozero) at (5.35,-.12);
  \coordinate (Oone) at (7.35,.12);

  \draw[singletonone] (4.58,-1)--(Ozero)--(4.58,1);
  \draw[singletontwo] (5.76,-1)--(Ozero)--(5.76,1);
  \draw[completion] (Ozero)--(4.88,1);
  \draw[completion] (Ozero)--(5.10,.64);
  \draw[completion] (Ozero)--(5.53,1);
  \node[focus,label={[black!70]right:$B_0$}] at (5.10,.64) {};

  \draw[singletonone] (6.94,-1)--(Oone)--(6.94,1);
  \draw[singletontwo] (8.12,-1)--(Oone)--(8.12,1);
  \draw[completion] (Oone)--(7.10,1);
  \draw[completion] (Oone)--(7.56,.70);
  \draw[completion] (Oone)--(7.82,1);
  \node[focus,label={[black!70]right:$B_1$}] at (7.56,.70) {};

  \node[focus,label={[black!70]above:$\widetilde A$}] at (6.08,-.62) {};
  \node[focus,label={[black!70]above:$\widetilde C$}] at (6.61,-.62) {};
  \fill[black] (Ozero) circle (1.45pt);
  \fill[black] (Oone) circle (1.45pt);
  \node[below right=0pt] at (Ozero) {$O_0$};
  \node[above left=0pt] at (Oone) {$O_1$};
  \node[anchor=south] at (6.35,1.22)
    {\textbf{(b)} two disjoint lifts in $\mathcal A^\circ$};
\end{scope}
\end{tikzpicture}

\end{adjustbox}
\caption{Cofactor placement.  The two full chords carry the singleton
factors.  Off-diagonal corrections exit the deleted boundary; the diagonal
support ends at a lift of $B$.  Crosswise gluing of the branch cut produces
two copies of the joint without adding intersections.}
\label{fig:cofactor-supports}
\end{figure}

\begin{proposition}\label{prop:proper-supports}
Modulo $I^N$, the cofactor placement is a proper locally finite wall support
system on $\Acal^\circ$.  Its only finite joints are the two lifts $O_0,O_1$
of $O$.  Distinct winding lifts do not meet on the annulus, even when their
developed images intersect in an affine plane.
\end{proposition}

\begin{proof}
Modulo $I^N$ there are finitely many deformation bidegrees.  After resonant
corrections have been grouped with the incoming chords, the compact polygon
contains a finite embedded chord-and-ray skeleton whose members meet only at
$O$.  In the normalization of Appendix~\ref{app:calculations}, the branch
segment joins $(-1/2,-1/2)$ to $(3/2,-1/2)$ and crosses the two negative
singleton chords at $(2/3,-1/2)$ and $(1/3,-1/2)$, away from $O$.
Crosswise gluing lifts each chord smoothly and produces two copies of the
skeleton, without creating another intersection.  Deleting the outer and
focus endpoints makes each support relatively closed.  A compact subset of
$\Acal^\circ$ stays a positive distance from these ends and therefore meets
finitely many supports.  Developed intersections of different deck lifts do
not descend to intersections of the supports themselves.
\end{proof}

Transport the ordered factorization of
Proposition~\ref{prop:singleton-completion} through the Jacobi path atlas.

\begin{theorem}[Open annular consistency]\label{thm:annular-consistency}
The transported factors on the supports of
Proposition~\ref{prop:proper-supports} define a deck-equivariant Jacobi wall
structure on $\Acal^\circ$.  It is consistent at every joint and
descends to
$\widehat\Xfrak_{\mathrm{ann}}^{\mathrm{root}}$.  Every path-ordered
transport along a compact transverse path modulo $I^N$ is defined on one
finite root-stack presentation and depends only on the contractible homotopy class of the
path relative to its endpoints.
\end{theorem}

\begin{proof}
At either joint $O_i$, the two full chords contribute the commutator in
\eqref{eq:singleton-commutator}, while the outgoing walls carry its unique
ordered inverse factorization.  Their product is the identity modulo $I^N$
for every $N$.  Changing a path chart conjugates this complete identity by a
Jacobi isomorphism.  Deck transport carries one wall skeleton to the
next, and Proposition~\ref{prop:proper-supports} excludes new joints between
different lifts.  This proves descent and contractible consistency.

A compact path crosses finitely many walls modulo $I^N$.  Theorem
\ref{thm:forced-category} places all of its root and wall transitions on one
finite stage.  Subdivision of a contractible homotopy into small joint loops
then gives homotopy invariance.
\end{proof}

The qualification ``contractible'' retains the global annular holonomy absent
from the boundary-charge cone.

\subsection{Hyperbolic holonomy and finite covers}

The matrix $K$ in \eqref{eq:intro-K} satisfies
\begin{equation}\label{eq:K-spectrum}
 \chi_K(T)=T^2-18T+1,
 \qquad \lambda_\pm=9\pm4\sqrt5.
\end{equation}
For $q_1=(2,1)$ and $q_2=(1,2)$, put $q_{i,n}=K^nq_i$.  Then
\begin{equation}\label{eq:Pell-recurrence}
 q_{i,n+2}=18q_{i,n+1}-q_{i,n},
 \qquad \det(q_{1,n},q_{2,n})=3.
\end{equation}
The core fixes their classes in $D$ but moves every nonzero exponent through
an infinite orbit.

More precisely, if a path $\gamma$ has linear holonomy $L_\gamma$, its
coefficient action in a developed chart has the form
\begin{equation}\label{eq:coefficient-holonomy}
 \Psi_\gamma(z^m)=c_\gamma(m)z^{L_\gamma m},
\end{equation}
where $c_\gamma$ is a multiplicative character with values in finite products
of path-indexed root units.  Its unique lift to the volume line has
multiplier
\begin{equation}\label{eq:line-holonomy}
 c_\gamma(\kappa)^{-1}z^{\kappa-L_\gamma\kappa}.
\end{equation}
Thus the annular holonomy is a semidirect action on exponents, roots, and the
volume line; a residual finite root phase alone does not describe it.

\begin{theorem}[Theta local system]\label{thm:theta-local-system}
No nonzero $m\in\widetilde M$ has finite $K$-orbit.  A single-valued theta
function on the annular quotient therefore cannot have a unique nonzero
leading monomial $z^m$.  Moreover, the formal orbit sum
\[
 \sum_{n\in\ZZ}\Hol(K)^n(\vartheta_m)
\]
does not converge in the $I$-adic completion.  The framed theta modules,
with the action \eqref{eq:coefficient-holonomy}--\eqref{eq:line-holonomy},
form a $K$-equivariant local system rather than a globally lattice-indexed
basis.
\end{theorem}

\begin{proof}
Neither eigenvalue in \eqref{eq:K-spectrum} is a root of unity, so
$K^r-1$ is invertible over $\RR$ for every $r\ne0$.  Thus no nonzero lattice
vector is periodic.  Holonomy takes a theta section with leading term $z^m$
to one with leading term $z^{Km}$, proving the first assertion.  The terms
of the orbit sum have distinct leading monomials, all of deformation order
zero.  Their orders do not tend to infinity, so the $I$-adic topology does
not sum them.
\end{proof}

Two finite covers retain the finite grading data.  Let
$p_{\mathrm{sec}}:\Acal^\circ_{\mathrm{sec}}\to\Acal^\circ$ be the
degree-two cover on which $D$ is constant.  The residual involution exchanges
$\alpha$ and $\beta$.  Put
\begin{equation}\label{eq:GKT-residue-group}
 M_{\GKT}=4\ZZ^2\subset M\subset\widetilde M,
 \qquad
 Q_{\GKT}=\widetilde M/M_{\GKT}\simeq(\ZZ/4)^2.
\end{equation}
Every peripheral monodromy acts on $Q_{\GKT}$ by the order-four matrix
\begin{equation}\label{eq:residue-matrix}
 P_{\mathrm{res}}=\begin{pmatrix}2&1\\3&0\end{pmatrix}\pmod4.
\end{equation}
The corresponding degree-four ordered-residue cover factors through
$p_{\mathrm{sec}}$.  Both covers preserve the pro-Kummer roots, volume line,
and wall structure.  Neither kills the hyperbolic exponent holonomy: the
first closed core lift to the ordered-residue cover acts by $K^2$.

On the sector cover, let $\delta_{\mathrm{deck}}$ denote the deck involution.  For a
$D$-graded deck-linearized object $\mathcal F$, define
\begin{equation}\label{eq:pushforwards}
 \Corr(\mathcal F)=p_*\mathcal F,
 \qquad
 \Sym(\mathcal F)_j
 =\left(\bigoplus_{\sigma(d)=j}p_*\mathcal F_d\right)^{\delta_{\mathrm{deck}}}.
\end{equation}
The first functor retains the two singleton labels and the deck operator;
the second takes the coordinate-symmetric summand.  Since the ground field
has characteristic zero, the idempotents
\[
 e_+=\frac{1+\delta_{\mathrm{deck}}}{2},
 \qquad e_-=\frac{1-\delta_{\mathrm{deck}}}{2}
\]
split the labelled object into symmetric and Prym parts.  Products are first
constructed equivariantly and only then restricted to $e_+$.

\subsection{Framed theta multiplication}

A \emph{path frame} consists of a lifted source chamber, an initial exponent
and generic endpoint, a compact region containing the broken lines under
consideration, a deformation order $N$, and a finite root-stack presentation supporting
the resulting finite diagram.  For such a frame define
\begin{equation}\label{eq:theta-sum}
 \vartheta_p=\sum_{\gamma\in\operatorname{BL}(p,Q)}
 a_\gamma z^{m(\gamma)}=z^p+O(I).
\end{equation}
Modulo $I^N$ the sum is finite: only finitely many walls meet the compact region,
and a contributing broken line has at most $N-1$ nontrivial bends.

\begin{theorem}[Finite-window theta multiplication]
\label{thm:theta-multiplication}
Fix finitely many framed theta sections, products, and associativity diagrams
modulo $I^N$.  After a finite enlargement of the exponent set and refinement
of the root-stack presentation, every requested product has a unique theta re-expansion.  These
products are commutative, associative, and equivariant under path and deck
transport.

If $B(z^m)=\vartheta_m$ is the finite change-of-basis matrix and
$b_{p,q}$ is the vector of monomial coefficients in
$\vartheta_p\vartheta_q$, then
\begin{equation}\label{eq:theta-reexpansion}
 B=1+U,
 \qquad B^{-1}=\sum_{j=0}^{N-1}(-U)^j\pmod{I^N},
 \qquad c_{p,q}=B^{-1}b_{p,q}.
\end{equation}
\end{theorem}

\begin{proof}
Proper local finiteness and the bend bound make every chosen theta sum and
ordinary product finite.  If re-expansion requires a new terminal exponent,
adjoin its theta sum and repeat.  Every correction raises $I$-order, so the
process stops after at most $N-1$ rounds.  Theorem
\ref{thm:forced-category} gives one finite root-stack presentation for the resulting
diagram.

The no-bend term makes $B$ unipotent, and the finite geometric series in
\eqref{eq:theta-reexpansion} is its two-sided inverse.  Re-expansion is
therefore unique.  Commutativity and associativity follow from ordinary
multiplication in the source chart.  Path and deck transports are algebra
isomorphisms and carry the complete change-of-basis diagram with them.
\end{proof}

The raw monomial vector $b_{p,q}$ and the theta structure constants
$c_{p,q}$ are different quantities.  This distinction will be essential in
Section~\ref{sec:primary-frobenius}, where Jacobian reduction and period
normalization supply two further changes of basis.
